% Chapter 15: The Kakeya Conjecture
\let\LocalMacro\relax

\chapter{Kakeya Conjecture}

\section{The Problem}

\subsection{Informal}

\textbf{Background.} A line segment has a direction---think of it as an arrow pointing from one end to the other. A region in the plane is just a collection of points forming a shape. Some shapes are small in the sense of having little area, yet they can still be strangely structured.

\textbf{Given.} You are allowed to draw a region in the plane that must contain at least one line segment pointing in every possible direction.

\textbf{Goal.} What is the smallest possible dimension of a region that contains a line segment pointing in every direction?

\textbf{Constraints.} The region must contain a line segment---not just a point---pointing in every direction from 0 to 180 degrees. The region can be disconnected and can have zero area. The question is about the intrinsic ``thickness'' of the set, measured by its dimension.

\textbf{Example.} Take a small disk in the plane. It contains line segments pointing in many directions, but not all---segments near the boundary point in limited directions. A Kakeya set is a shape that contains a segment in every direction, yet (surprisingly) can have zero area. The conjecture says that even though the area can be zero, the set must still be ``thick'' in the sense of having the full dimension of the plane.

\subsection{Formal}
\begin{conjecture}[Kakeya Conjecture]
Any Kakeya set in $\mathbb{R}^n$ must have Hausdorff dimension and Minkowski dimension equal to $n$.
\end{conjecture}

Despite its geometric statement, the Kakeya conjecture is deeply connected to harmonic analysis, specifically to wave equations, Fourier multipliers, and the restriction conjecture.

\subsection{Historical Context}
The Kakeya conjecture is one of the most famous unsolved problems in analysis. While it is settled for $n=2$, it remains open in higher dimensions ($n \geq 3$), representing a key barrier to understanding multi-directional wave propagation. Nets Katz of Rice University called the 3D progress the solution of the ``holy grail'' problem in harmonic analysis.

\subsection{What Was Known}
Besicovitch proved in 1919 that Kakeya sets can have zero area, which was already shocking. The question then became: how small can their dimension be? In two dimensions, Davies proved in 1971 that Kakeya sets must have dimension 2. In higher dimensions, the conjecture remained wide open. Finite field analogues were proved by Dvir in 2008, but the real case resisted all attacks.

\subsection{Why It Matters}
The Kakeya conjecture is about the size of the smallest region that can contain a needle pointing in every direction. It connects to harmonic analysis (the study of waves), partial differential equations (the study of fluid flow), and even quantum physics. Understanding how thick Kakeya sets must be has implications for the behavior of solutions to the wave equation.

\subsection{Simplified Version}
The core idea is that containing a line segment in every direction forces a set to be ``thick'' in a precise geometric sense. In two dimensions, this is already proved: any Kakeya set in $\mathbb{R}^2$ must have Hausdorff dimension 2, meaning it is as large as a full-dimensional region even though it can have zero area. The conjecture asserts that this holds in all dimensions---Kakeya sets in $\mathbb{R}^n$ must have dimension $n$. The finite field analogue, proved by Dvir in 2008 using the polynomial method, was a crucial stepping stone toward the real case.

\section{Mathematical Theory}

\subsection{Prerequisites}
This chapter assumes familiarity with:
\begin{itemize}
\item Basic real analysis (Lebesgue measure, open/closed sets, limits)
\item Geometric measure theory (Hausdorff dimension)
\item Finite fields and polynomials (roots, finite vector spaces)
\item Harmonic analysis (Fourier transform, restriction operators) --- helpful but not essential
\end{itemize}

\subsection{Besicovitch's Construction}
Besicovitch's proof uses a geometric decomposition called ``panning.'' By dividing a triangle into smaller triangles, shifting them, and overlapping them, one can concentrate the area into arbitrarily small regions while preserving segments in all directions.

\subsection{Finite Field Analogue}
In 1999, Thomas Wolff proposed studying the Kakeya conjecture in the simplified setting of vector spaces over finite fields $\mathbb{F}_q^n$ to bypass real-variable complications \cite{wolff1999}.

\begin{definition}[Finite Field Kakeya Set]
A subset $K \subset \mathbb{F}_q^n$ is a Kakeya set if for every direction $y \in \mathbb{F}_q^n \setminus \{0\}$, there exists $x \in \mathbb{F}_q^n$ such that $\{x + t y : t \in \mathbb{F}_q\} \subset K$.
\end{definition}

Wolff conjectured that any Kakeya set in $\mathbb{F}_q^n$ satisfies $|K| \geq C_n q^n$.

\section{The Solution}

\subsection{The Big Idea}
A Kakeya set in $\mathbb{R}^n$ contains a line segment in every direction. Besicovitch showed in 1919 that it can have zero volume. Davies proved the two-dimensional case in 1971, and Wang and Zahl established the full Hausdorff and Minkowski dimension bound in $\mathbb{R}^3$ in 2025. Their proof uses multiscale geometric estimates and incidence theory rather than merely the earlier $2.5+\epsilon$ lower bound.

\subsection{What Was Stopping Progress}
In two dimensions, Davies proved in 1971 that Kakeya sets must have full dimension 2. In three dimensions, Fourier-analytic methods---restriction theory, decoupling, endpoint estimates---had long hit a ceiling around 2.5 before the Wang--Zahl work.

The polynomial method that Dvir used so brilliantly for the finite field case (2008) had no direct translation to the real setting. Over a finite field $\mathbb{F}_q$, a polynomial of degree less than $q$ that vanishes on a line must vanish everywhere---this is a finiteness fact: a nonzero univariate polynomial of degree $d$ has at most $d$ roots, so if $d < q$ and the polynomial vanishes on all $q$ points of a line, it must be the zero polynomial. But over $\mathbb{R}$, lines have infinitely many points and there is no degree bound---a nonzero polynomial can vanish on an entire line without being zero. This means Dvir's argument, which directly constructed a low-degree polynomial vanishing on a Kakeya set, cannot work over $\mathbb{R}$. The challenge was to find a way to use polynomials without requiring them to vanish on entire lines.

\subsection{The Breakthrough}
Wang and Zahl combined multiscale estimates for tube configurations with incidence geometry to prove that Kakeya sets in $\mathbb{R}^3$ have Hausdorff and Minkowski dimension 3. The full conjecture in dimensions $n\geq4$ remains open.

\subsection{How It Works}

\subsection{Dvir's Finite Field Proof (2008)}
In 2008, Zeev Dvir proved Wolff's finite field conjecture in a paper that was only four pages long, using the \emph{polynomial method} \cite{dvir2009}:

\begin{theorem}[Dvir, 2008 \cite{dvir2009}]
Any Kakeya set $K \subset \mathbb{F}_q^n$ satisfies:
\begin{equation}
|K| \geq \frac{1}{n!} q^n
\end{equation}
\end{theorem}

The proof is elegant: if $|K| < \binom{n+q-1}{n}$, there exists a non-zero polynomial of degree $<q$ vanishing on $K$. By complete induction along line directions, the polynomial's highest degree homogeneous part must vanish on all of $\mathbb{F}_q^n$, which is impossible.

\subsection{The 3D Real Breakthrough: Wang--Zahl (2025)}

The polynomial method had been spectacularly successful for the finite field Kakeya problem, but extending it to the real setting $\mathbb{R}^n$ had resisted all attempts for decades. The breakthrough came from Hong Wang (NYU) and Joshua Zahl (UBC).

\begin{historicalbox}
Hong Wang (born 1992, China) brought deep expertise in polynomial partitioning and geometric measure theory. Joshua Zahl (UBC) contributed experience in incidence geometry and harmonic analysis. Their collaboration spanned several years and required new ideas at each stage.
\end{historicalbox}

\begin{theorem}[Wang--Zahl, 2025 \cite{wangzahl2025}]
Any Kakeya set $K \subset \mathbb{R}^3$ has Hausdorff and Minkowski dimension $3$.
\end{theorem}

Their proof combined three major innovations:

\begin{enumerate}
\item \textbf{Polynomial partitioning}: A technique from algebraic geometry (introduced by Guth and Katz in 2010) that decomposes space using the zero sets of polynomials. Given a Kakeya set $K \subset \mathbb{R}^3$, construct a polynomial $P$ that vanishes on a carefully chosen subset of the lines in $K$. The zero set of $P$ decomposes space into cells, each containing relatively few lines.
\item \textbf{Linear programming bounds}: Optimizing over families of line configurations to extract dimensional lower bounds. This technique converts geometric statements about Kakeya sets into optimization problems that can be solved systematically.
\item \textbf{New incidence geometry}: Novel bounds on the number of incidences between points and lines in three dimensions, extending the work of Guth and Katz.
\end{enumerate}

\begin{highlightbox}
The Wang--Zahl result settles the Kakeya dimension conjecture in three dimensions. The corresponding Euclidean conjecture in dimensions $n\geq4$ remains open.
\end{highlightbox}

\subsection{Hong Wang and the Fields Medal}

In 2026, Hong Wang was awarded the Fields Medal at the International Congress of Mathematicians in Philadelphia, becoming the third woman to receive the prize (after Maryam Mirzakhani in 2014 and Maryna Viazovska in 2022). Her work demonstrated that geometric and combinatorial techniques can attack real-variable problems that had resisted Fourier-analytic methods for decades. Kakeya-type configurations remain central to the restriction conjecture and to estimates for wave propagation.

\subsection{After the Proof}

The Kakeya conjecture is settled in dimensions 2 and 3, and in the finite-field model. In $\mathbb{R}^3$, Wang and Zahl proved in 2025 that Kakeya sets have Hausdorff and Minkowski dimension 3. The general Euclidean problem in dimensions $n\geq4$ remains open.

The Wang--Zahl breakthrough broke a barrier that had stood for decades. Their multiscale geometric and incidence arguments establish the full dimension bound in $\mathbb{R}^3$, opening a new front of attack on the higher-dimensional problem.

The full conjecture---that Kakeya sets in $\mathbb{R}^n$ have Hausdorff dimension $n$ for all $n$---remains open in dimensions $n\geq4$. The three-dimensional proof does not automatically generalize, and the higher-dimensional problem is deeply connected to the restriction conjecture and to bounds for wave equations.

\section{Impact}
\begin{itemize}
\item \textbf{Harmonic Analysis}: Directly bounds the restriction operator for the Fourier transform, influencing local smoothing estimates for wave equations. The Wang--Zahl bound limits how concentrated energy can become during wave propagation.
\item \textbf{Additive Combinatorics}: The finite field polynomial method has been adapted to prove the joints conjecture and other incidence geometry problems. The incidence geometry techniques have applications to sum-product estimates.
\item \textbf{Partial Differential Equations}: Provides bounds on the size of singularity sets for dispersive wave propagation.
\item \textbf{Geometric Measure Theory}: The polynomial method developed for Kakeya has been applied to Falconer's distance set problem and other classical questions.
\end{itemize}

\section{Exercises}

\begin{exercise}[Basic]
Show that a Kakeya set in $\mathbb{R}^2$ containing segments in all directions can be constructed with measure 0.
\end{exercise}

\begin{exercise}[Intermediate]
Write down the details of Zeev Dvir's finite field Kakeya set proof. Show why the degree of the vanishing polynomial must be less than $q$.
\end{exercise}

\begin{exercise}[Challenge]
Explain why Dvir's algebraic proof fails in the setting of real numbers $\mathbb{R}^n$. What is the main difference between lines in $\mathbb{F}_q^n$ and $\mathbb{R}^n$?
\end{exercise}

\begin{exercise}[Computational]
Write a Python/SageMath script to explore a concept from this chapter. See the appendix for starter code.
\end{exercise}

\section{Timeline}

\begin{tabularx}{\linewidth}{lX}
\toprule
\textbf{Year} & \textbf{Event} \\ \midrule
1917 & Kakeya poses the needle turning problem \cite{kakeya1917} \\
1919 & Besicovitch constructs a Kakeya set of measure zero \cite{besicovitch1919} \\
1971 & C\'ordoba proves the Kakeya maximal function conjecture for $n=2$ \\
1999 & Wolff formulates the finite field Kakeya conjecture \cite{wolff1999} \\
2003 & Wolff improves the lower bound on Kakeya dimension in $\mathbb{R}^3$ \\
2008 & Dvir proves the finite field Kakeya conjecture using the polynomial method \cite{dvir2009} \\
2025 & Wang and Zahl prove the three-dimensional Kakeya dimension conjecture \cite{wangzahl2025} \\
2026 & Hong Wang awarded the Fields Medal at ICM Philadelphia \\
\bottomrule
\end{tabularx}

\section{Citations}

\begin{enumerate}
\item Kakeya, S. (1917). ``Some problems on maximum and minimum regarding ovals.'' Tohoku Mathematical Journal.
\item Besicovitch, A. S. (1919). ``Sur deux questions de geometrie.'' Bulletin de l'Academie des Sciences de Russie.
\item Wolff, T. (1999). ``Recent work on restriction and Kakeya conjectures.'' IAS Lectures.
\item Dvir, Z. (2009). ``On the size of Kakeya sets in finite fields.'' Journal of the AMS, 23(4), 1107--1119.
\item Guth, L., \& Katz, N. (2015). ``On the Erd\H{o}s distinct distances problem in the plane.'' Annals of Mathematics, 181(1), 155--190.
\item Wang, H., \& Zahl, J. (2025). ``The Kakeya conjecture in three dimensions.'' arXiv:2502.17655.
\item International Mathematical Union. (2026). ``Fields Medal citation: Hong Wang.'' ICM 2026, Philadelphia.
\end{enumerate}
